% !PRINTED 272
\weberpage{272}
%<<CONT
ist, und jeder Theiler dieser Zahlen tritt auch unter den Graden der cyklischen Gruppen auf. Der Index eines cyklischen Theilers kann niemals kleiner sein als $1/2\,(p^2-1)$, $p(p+1)$, $p(p-1)$.

Wenn der Grad eines Theilers $G$ von $L_p$ durch $p$ theilbar ist, so enthält er eine cyklische Gruppe vom Grade $p$, und wir können nach \S.\,68 die Gruppe $G$ so transformiren, dass diese cyklische Gruppe aus den Substitutionen
\[
T = \begin{pmatrix}1, & t\\ 0, & 1\end{pmatrix}, \quad t = 0, 1, 2, \ldots, p-1
\]
besteht. Wenn nun $G$ noch eine Substitution
\[
A = \begin{pmatrix}a, & b\\ c, & d\end{pmatrix}
\]
enthält, in der $c$ von Null verschieden ist, so kommt darin auch
\weq{1}{\begin{pmatrix}1, & -ac^{-1}\\ 0, & 1\end{pmatrix}\begin{pmatrix}a, & b\\ c, & d\end{pmatrix}\begin{pmatrix}1, & -dc^{-1}\\ 0, & 1\end{pmatrix} = \begin{pmatrix}0, & -c^{-1}\\ c, & 0\end{pmatrix}}
vor. Daraus folgt weiter, dass auch
\[
\begin{pmatrix}0, & -c^{-1}\\ c, & 0\end{pmatrix}\begin{pmatrix}1, & t\\ 0, & 1\end{pmatrix}\begin{pmatrix}0, & -c^{-1}\\ c, & 0\end{pmatrix} = \begin{pmatrix}1, & 0\\ -c^2 t, & 1\end{pmatrix},
\]
und mithin jede Substitution
\[
U = \begin{pmatrix}1, & 0\\ u, & 1\end{pmatrix}, \quad u = 0, 1, 2, \ldots, p-1
\]
vorkommt. Setzt man $U$ für $u = -1$ an Stelle von $A$ in die Formel (1), so folgt, dass $G$ auch die Substitution $\begin{pmatrix}0, & 1\\ -1, & 0\end{pmatrix}$ enthält, die mit $U$ zusammen ein System erzeugender Elemente von $L_p$ giebt (\S.\,65), so dass also $G$ mit $L_p$ identisch ist. Es folgt hieraus, dass $G$, wenn es nicht mit $L_p$ identisch sein soll, nur Substitutionen von der Form
\[
A = \begin{pmatrix}a, & b\\ 0, & a^{-1}\end{pmatrix}
\]
enthalten kann. Umgekehrt bilden alle diese Substitutionen eine Gruppe vom Grade $1/2\,p(p-1)$, also einen Theiler vom Index $p+1$.

Es sind darunter noch gewisse Theiler von grösserem Index enthalten, die man erhält, wenn man für $a$ nicht alle Werthe nimmt, sondern alle Werthe von der Form $a^s$, wenn $s$ ein Theiler von $1/2\,(p-1)$ ist. Der Grad einer solchen Gruppe ist $1/2\,p(p-1) : s$.
